% Sections 7.1-7.5, from lectures/L07/appendix/a_timers.md (A.1-A.5).

\section{Timers: a counter with a target value}
\label{c7:sec:concept}

A \textbf{timer} answers a different question from the counter in \secref{c6:sec:counter}. A counter
answers ``what is the count right now?''; a timer answers ``has a certain amount of time passed?''

It does so with three parts: an internal counter invisible from outside, a target value
\code{TICK_COUNT} fixed per instance, and a comparison which, when true, clears the counter and
pulses the output flag \code{timeout} for one clock cycle.

Because the system clock ticks at a known, fixed frequency, counting ticks is the same thing as
measuring elapsed time. On the DE0-CV that is 50,000,000 ticks per second:

\[ \mathit{TICK\_COUNT} = \mathit{seconds} \times 50\,000\,000 \]

\begin{itemize}
  \item \code{TICK_COUNT = 50_000_000}: timeout once per second.
  \item \code{TICK_COUNT = 25_000_000}: twice per second.
  \item \code{TICK_COUNT = 5_000_000}: ten times per second.
\end{itemize}

Strictly speaking the counter runs \code{0} through \code{TICK_COUNT} inclusive, so a full period is
\code{TICK_COUNT + 1} cycles. That extra tick is between 0.02 and 0.2~ppm at these values, far below
the board oscillator's own tolerance, so the round numbers above are used for clarity; for an exact
cycle count, use $\mathit{TICK\_COUNT} = \mathit{seconds} \times 50\,000\,000 - 1$.

This is the opposite of a blocking \code{_delay_ms()}. A blocking delay stops the whole program to
wait; a timer runs permanently in parallel with everything else and simply raises a flag, and
nothing in the design is ever blocked waiting for it.

% ------------------------------------------------------------------------------------------
\section{Building a timer by hand in CircuitVerse}
\label{c7:sec:byhand}

Before any VHDL is written: build a timer as a gate network, exactly as you did for \chapref{c1:ch}
and \ref{c2:ch}'s combinational networks and \chapref{c3:ch}'s flip-flops. A timer is small enough
to draw in its entirety, and drawing it is what makes \secref{c7:sec:module}'s VHDL read as a
description of something you already understand rather than a new idiom to learn by heart.

Build it in a size you can follow: a \textbf{4-bit counter} with \code{TICK_COUNT = 10}, which is
\code{1010}. The real design counts to 50,000,000; nothing in the structure changes, only the width.

\textbf{Start from the counter you have already built} in \secref{c6:sec:counter}: a 4-bit register
\code{Q0} to \code{Q3} sharing a clock, an adder and a constant \code{1} with the sum fed back to the
\code{D} inputs, and nothing more. If you did not keep it, it is three elements and about five
minutes to redraw. Confirm that it still counts and wraps before you add anything, so that you know
the counter underneath is sound when the timer later misbehaves.

Everything here adds \textbf{two things on top of that}: a comparison saying ``the count has
arrived'', and a clear making it repeat. That is the whole difference between a counter and a timer.

\paragraph{Add the comparison}
You want \code{timeout} high exactly when the four bits spell \code{1010}, and the clean way to
build ``are these two numbers equal?'' is one \textbf{XNOR gate per bit}. An XNOR gives \code{1}
when its inputs are \emph{the same}, which is \secref{c1:sec:gates}'s XNOR column read as an equality
test. Wire in four, each comparing one counter output against the corresponding target bit
(\code{Q3} against \code{1}, \code{Q2} against \code{0}, \code{Q1} against \code{1}, \code{Q0}
against \code{0}), and drive the target side from a constant in CircuitVerse rather than from a wire
you might later mistake for a signal. Then combine the four outputs with AND; that single AND gate
is \code{timeout}. With $\odot$ for XNOR, the one operator \secref{c1:sec:boolean} gave no symbol:

\[ \mathit{timeout} = (Q_3 \odot 1) \cdot (Q_2 \odot 0) \cdot (Q_1 \odot 1) \cdot (Q_0 \odot 0) \]

All four bits have to match simultaneously, which is what ``the counter has reached \code{1010}''
means. This is a \textbf{general equality comparator}, worth recognising as a building block of its
own: the same four XNORs plus an AND compare against any 4-bit value, and only the constants change.

\paragraph{Add the clear}
So that the timer restarts instead of carrying on: feed \code{timeout} back to force the counter to
\code{0000} at the next edge, by gating every flip-flop's \code{D} input so that \code{timeout = 1}
takes precedence over the adder's output. Without it the counter carries on to \code{1111} and wraps
by itself, so \code{timeout} fires every sixteenth edge instead of every eleventh. Both are
periodic; only one has the period you asked for.

\paragraph{Then confirm, one clock edge at a time}
\begin{itemize}
  \item The counter goes from \code{0000} through \code{1010}, and \code{timeout} is low for every
    count value except \code{1010}. Keep an eye on the four XNOR outputs as it counts and note how
    often three are high at once. Three bits matching is no hit; it is the AND gate insisting on all
    four.
  \item At the edge after \code{timeout} fires, the counter is back at \code{0000} and
    \code{timeout} is low, so it is high for exactly one clock cycle, never two.
  \item The whole cycle is therefore eleven edges, count values \code{0} through \code{10}, repeated
    forever. That is \secref{c7:sec:concept}'s \code{TICK_COUNT + 1} off-by-one, and this is the
    book's easiest place to see it, since you can count the edges by eye.
\end{itemize}

\lecfig[0.5]{lectures/L07/appendix/images/timer_circuit.png}{The comparison logic: four XNORs
against \code{1010}, ANDed into \code{timeout}.}{c7:fig:timercircuit}

The drawing is cropped to the part this section adds. It labels the counter's bits
\code{counter[3]} down to \code{counter[0]} where the text above calls them \code{Q3} to \code{Q0},
and the increment feeding them is drawn as gate-level carry logic rather than the single adder block
\secref{c6:sec:counter} used, since that is what CircuitVerse gives you when you build it one bit at
a time. The four XNOR gates, their constants \code{1}, \code{0}, \code{1}, \code{0}, and the AND
combining them are all that is new here.

Two things are worth noticing before the VHDL, since both are what the language buys you:
\begin{itemize}
  \item \textbf{The comparator is hardwired to one target value.} Changing \code{TICK_COUNT} from
    \code{10} to \code{12} means going back into the drawing and flipping two constants. It is a
    small change, but it is a change to the \emph{circuit}, and every instance needs its own copy.
    In \secref{c7:sec:module} the target value becomes a generic.
  \item \textbf{The width is hardwired too.} Counting to 50,000,000 takes 26 flip-flops and a 26-bit
    comparator, 26 XNORs over 52 inputs: the same circuit, an unreasonable drawing. In
    \secref{c7:sec:module} it is the same six lines of VHDL either way.
\end{itemize}

That is the trade-off this book makes throughout: draw the circuit once, at a size you can see, to
know what the tool builds for you, and then let the tool build it at a size you could not have
drawn.

% ------------------------------------------------------------------------------------------
\section{The \code{timer} module in VHDL}
\label{c7:sec:module}

The same circuit as a reusable module. It is built live during the session, and
Exercise~\ref{c7:ex:timer} asks you to write it yourself afterwards, so it is printed here in full
rather than handed out as a file:

\begin{vhdlcode}
entity timer is
    generic(TICK_COUNT: natural := 50_000_000);
    port(clock, reset_s2_n, enable: in std_logic;
         timeout                  : out std_logic);
end entity;

architecture behaviour of timer is
-- Internal tick counter.
signal counter: natural range 0 to TICK_COUNT;
begin
    process(clock, reset_s2_n) is
    begin
        if (reset_s2_n = '0') then
            timeout <= '0';
            counter <= 0;
        elsif (rising_edge(clock)) then
            timeout <= '0';
            if (enable = '1') then
                if (TICK_COUNT > counter) then
                    counter <= counter + 1;
                else
                    timeout <= '1';
                    counter <= 0;
                end if;
            end if;
        end if;
    end process;
end architecture;
\end{vhdlcode}

\begin{itemize}
  \item \code{TICK_COUNT} is a \textbf{generic}, not a constant baked into the architecture, so the
    same entity serves any frequency by being instantiated with a different \code{generic map}. The
    internal counter's range is written in terms of it, so the register is sized for the target value
    rather than for whatever an unconstrained \code{natural} would give, the same point
    \secref{c6:sec:counter} makes about \code{natural range 0 to 15}. A generic can size a signal
    precisely because it is fixed before synthesis runs.
  \item \code{reset_s2_n} is the already \emph{synchronised} reset signal, never the raw
    asynchronous \code{reset_n}. A timer never touches a raw asynchronous input directly.
  \item \code{enable} \textbf{freezes} the timer rather than clearing it. While it is \code{'0'},
    \code{counter} holds its value and \code{timeout} stays low, and a disabled timer carries on
    where it left off rather than restarting. \Chapref{c8:ch}'s \code{fsm_led} depends on that.
  \item \code{timeout} is a one-cycle pulse. The \code{else} branch setting it is reached at only
    one edge: the edge where \code{counter} already stands at \code{TICK_COUNT} and is cleared back
    to \code{0}. At the following edge the unconditional \code{timeout <= '0';} runs with nothing
    after it to overwrite it, so the pulse is exactly one cycle wide.
\end{itemize}

Compare this with \secref{c7:sec:byhand}'s gate network. The logic is the same, and
\secref{c7:sec:byhand} already named what the generic buys. One difference it did not name:
\secref{c7:sec:byhand}'s \code{timeout} is a combinational decode, high during the cycle where the
count shows \code{1010}, while this one is registered, high one cycle later, during the cycle where
\code{counter} shows \code{0}. Same period, different phase, and different glitch behaviour, which
is the distinction \chapref{c8:ch}'s discussion of Moore versus Mealy rests on.

% ------------------------------------------------------------------------------------------
\section{Composing a complete timer circuit}
\label{c7:sec:compose}

A timer on its own only pulses \code{timeout}; something has to react to that pulse, and a button
has to control the timer. With \chapref{c4:ch}'s two-flip-flop synchroniser reused, a complete
button-controlled, timer-driven LED circuit looks like this:

\lecfig[1.0]{lectures/L07/appendix/images/led_toggle_timer.png}{The whole circuit at once:
synchronised button, timer and toggled LED. The details are legible in the two close-ups
below.}{c7:fig:whole}

Read it as a floor plan rather than as a wiring diagram. It is a wide circuit, so at page size its
labels come out small; what it is for is showing how few parts there are and roughly where each
sits. The two close-ups below are where the signal names are readable, and together they cover
everything in it.

From left to right:
\begin{itemize}
  \item \textbf{Synchronisation and edge detection.} \code{button_n} passes through two flip-flops
    for metastability protection, and a third then lets the circuit compare ``now'' against
    ``previous'' and detect a falling edge, giving a one-cycle \code{button_edge_s2} pulse.
  \item \textbf{The timer} (\secref{c7:sec:module}), enabled or disabled by \code{timer_enabled},
    which in turn is toggled by \code{button_edge_s2}.
  \item \textbf{Toggling the LED.} The timer's \code{timeout} pulse flips the LED's state every time
    it fires, and the LED is forced off as soon as the timer is disabled or the circuit is reset.
\end{itemize}

\lecfig[0.98]{lectures/L07/appendix/images/button_sync_and_flank_detection.png}{Synchroniser and
falling-edge detector.}{c7:fig:buttonsync}

\lecfig[0.6]{lectures/L07/appendix/images/led_toggle.png}{The LED toggle logic, driven by the
timer's \code{timeout} pulse.}{c7:fig:ledtoggle}

Two patterns matter here, independently of this circuit:
\begin{itemize}
  \item \textbf{Every asynchronous input is synchronised before it touches any other logic}, and
    every enable or toggle signal derived from a button is a synchronous signal from that point on.
    You reuse that unchanged in \secref{c7:sec:walking} and again for \chapref{c8:ch}'s state
    machines.
  \item \textbf{One clock, and everything else is an enable.} \code{timeout} clocks the LED by being
    a one-cycle \emph{enable} for a flip-flop running on the system clock, not by being wired into a
    clock input. It is tempting to run a counter or timer output into a clock port to ``slow things
    down''; do not. A design with one clock is a design whose timing the tool can analyse, and this
    is the single most common mistake a programmer makes when a baud generator has to be written.
\end{itemize}

% ------------------------------------------------------------------------------------------
\section{\code{walking_led}: a shift register with a walking LED}
\label{c7:sec:walking}

\code{walking_led} is the first design composing three chapters' building blocks into a circuit that
can be demonstrated in hardware: a single lit LED walking along a row, one position per timer tick,
started and stopped with a button. It is built live during the session, and
Exercise~\ref{c7:ex:walking} asks you to write it yourself afterwards, so it is printed here in full
rather than handed out as a file.

It writes almost no new logic, reusing \code{reset_sync} and \code{button_sync} unchanged
(\secref{c4:sec:resetsync} and \ref{c4:sec:buttonsync}) to synchronise and edge-detect the button,
\code{timer} (\secref{c7:sec:module}) to clock the speed, and the shift register from
\secref{c6:sec:shift}--\ref{c6:sec:piso} for the walking itself. That is what the last four chapters
have accumulated: a small set of composable modules.

You write all three yourself: \code{reset_sync} and \code{button_sync} in
Exercise~\ref{c4:ex:resetsync} and \ref{c4:ex:buttonsync}, \code{timer} in
Exercise~\ref{c7:ex:timer} in this chapter. Copy your own into the exercise directory before you
build this one. Reusing a module presupposes having one, and this is the first design that asks you
for the ones you have already built.

\begin{vhdlcode}
entity walking_led is
    generic(LED_COUNT : natural := 8;
            TICK_COUNT: natural := 25_000_000);
    port(clock, reset_n, button_n: in std_logic;
         led                     : out std_logic_vector(LED_COUNT-1 downto 0));
end entity;

architecture behaviour of walking_led is
signal reset_s2_n, button_edge_s2  : std_logic;
signal button_n_v, button_edge_s2_v: std_logic_vector(0 downto 0);
signal shift_enable, shift_tick    : std_logic;
signal shift_reg: std_logic_vector(LED_COUNT-1 downto 0);

begin
    led <= shift_reg;

    button_n_v(0)  <= button_n;
    button_edge_s2 <= button_edge_s2_v(0);

    reset_sync1: entity work.reset_sync
        port map(clock, reset_n, reset_s2_n);

    button_sync1: entity work.button_sync
        generic map(1)
        port map(clock, reset_s2_n, button_n_v, button_edge_s2_v);

    timer1: entity work.timer
        generic map(TICK_COUNT)
        port map(clock, reset_s2_n, shift_enable, shift_tick);

    SHIFT_ENABLE_PROCESS: process(clock, reset_s2_n) is
    begin
        if (reset_s2_n = '0') then
            shift_enable <= '0';
        elsif (rising_edge(clock)) then
            if (button_edge_s2 = '1') then
                shift_enable <= not shift_enable;
            end if;
        end if;
    end process;

    SHIFT_PROCESS: process(clock, reset_s2_n) is
    begin
        if (reset_s2_n = '0') then
            shift_reg <= (0 => '1', others => '0');
        elsif (rising_edge(clock)) then
            if (shift_tick = '1') then
                shift_reg <= shift_reg(LED_COUNT-2 downto 0) & shift_reg(LED_COUNT-1);
            end if;
        end if;
    end process;
end architecture;
\end{vhdlcode}

Design choices worth noting:
\begin{itemize}
  \item \code{button_n_v} and \code{button_edge_s2_v} only bridge a width difference:
    \code{walking_led}'s button signals are single bits while \code{button_sync}'s ports are vectors
    of width \code{COUNT}. The two concurrent assignments connect the scalar to element \code{0} and
    back. They cost no hardware.
  \item This is a PISO register (\secref{c6:sec:piso}) in its simplest form: parallel-loaded exactly
    once, at reset, with a single lit bit at position 0, and shifting only thereafter. There is no
    \code{load} at run time, since this design never loads a new pattern.
  \item It is a \textbf{circular} shift register. The shift is \secref{c6:sec:shift}'s expression
    unchanged, so the lit bit walks towards the MSB; the only difference from an ordinary PISO is
    what feeds bit 0, which is the bit falling off the top rather than a \code{serial_in} from
    somewhere else. A one-shot register that ``shifts a pattern out'' becomes an endlessly repeating
    pattern generator for free, just by recirculating its own output.
  \item \code{shift_enable} is \secref{c7:sec:compose}'s toggle flip-flop reused verbatim: the
    button's synchronised edge pulse flips a single bit, and that bit becomes the timer's
    \code{enable}.
  \item \code{shift_tick} is the timer's \code{timeout} pulse, reused as the shift register's clock
    enable instead of driving an LED toggle. Same building block, a different consumer, and exactly
    the ``one clock, everything else is an enable'' rule from \secref{c7:sec:compose}.
  \item Both \code{LED_COUNT} and \code{TICK_COUNT} are generics, so
    \code{generic map(4, 25_000_000)} gives a row of 4 LEDs and \code{generic map(8, 5_000_000)} a
    faster 100~ms step, without touching the module.
\end{itemize}

On the board, demonstrated during the session, the ports are assigned via the Pin Planner
(appendix~\ref{app:quartus}): \code{clock} to the 50 MHz oscillator, \code{reset_n} and
\code{button_n} to two push buttons, and \code{led} to a row of LEDs on the board. Press the button
once to set the bit moving, once more to freeze it.

You can see the same behaviour without a board by running the reference testbench, which overrides
both generics so that a full round takes a handful of clock cycles:

\begin{shell}
cd lectures/L07/exercises/walking_led
cp ../../../L04/exercises/reset_sync/reset_sync.vhd .   # the three you wrote yourself
cp ../../../L04/exercises/button_sync/button_sync.vhd .
cp ../timer/timer.vhd .
ghdl -a --std=93 reset_sync.vhd button_sync.vhd timer.vhd \
                 walking_led.vhd walking_led_tb.vhd
ghdl -e --std=93 walking_led_tb
ghdl -r --std=93 walking_led_tb --assert-level=error --stop-time=10ms
\end{shell}
